% Reconstruction of the fully audited boundedness and fixed-rank arguments.
\begin{proposition}[Proved: Bounded minimal extensions of periodic reflexive sheaves]
\label{D:global-periodic-reflexive-boundedness}
Let $S$ be a normal proper surface over an algebraically closed field
of characteristic $p>0$. Fix a projective resolution $\pi:Y\to S$,
with smooth projective $Y$, an ample divisor $H$, and an effective
exceptional divisor $Z$ with $Z\cdot E_i<0$ on every exceptional
component. For reflexive sheaves $E$ of rank $r$ satisfying
$F_S^{[p^f]*}E\simeq E$ for some $f\ge1$, their minimal extensions
$F$ to $Y$ have the following bounds, with constants independent of
$r$ and $f$:
\begin{enumerate}
\item their restrictions to the exceptional curves have fixed upper
and lower slope windows;
\item $|\chi(Y,F)|\le Cr$, and $c_1(F)$ is numerically an exceptional
linear combination whose coefficients have absolute values at most $Cr$;
\item $\mu_{\max,H}(F)$ and $\mu_{\max,H}(F^\vee)$ are bounded by a
fixed constant, and the analogous bounds hold after restriction to a
fixed sufficiently ample smooth curve;
\item for each fixed $r$, these extensions form a bounded family.
\end{enumerate}
Ordinary isomorphism of their underlying reflexive sheaves, and the
one-step relation $F_S^{[*]}E\simeq E'$, are constructible relations on
one finite-type parameter space for each fixed rank.
\end{proposition}
\begin{proof}
We use minimal extensions in the sense of
\citet[Definition 1.2 and Lemma 1.3]{C_Langer2026}. On an exceptional
curve $E_i$ of genus $g_i$, their slopes lie in a fixed interval, which
we may take to have endpoints $g_i-1$ and $g_i-1-E_i^2$. The important
point is that the endpoints do not depend on the rank or on the
Frobenius period.

We record the uniform extension calculation. Filter $\mathcal O_Z$ by
a fixed finite sequence whose quotients are line-bundle twists on the
reduced exceptional components. For two bundles with the stated slope
windows, Serre duality and the elementary implication
\[
 \mu_{\min}(A)>\mu_{\max}(B)\quad\Longrightarrow\quad
 \operatorname{Hom}(A,B)=0
\]
give, for all sufficiently large fixed $n$,
\[
 H^0\bigl(Z,\operatorname{Hom}(F,G)\otimes\mathcal O_Z(nZ)\bigr)=0,
 \qquad
 H^1\bigl(Z,\operatorname{Hom}(F,G)\otimes\mathcal O_Z(-nZ)\bigr)=0.
\]
We apply the inequalities to the individual component layers; we do
not use an incorrect assertion that tensor products preserve
semistability in characteristic $p$. The first vanishing says that
every map away from the exceptional set has a fixed pole bound $BZ$.
Indeed such a map first extends with some finite pole order by Cartier
localization, and the exact sequences for successive $Z$-layers lower
that order to $B$. The second vanishing gives a fixed finite stage
which determines formal maps and bundles. Existence of lifts of bundles
uses the vanishing of $H^2$ on the one-dimensional exceptional support;
uniqueness and lifting of prescribed maps use the displayed $H^1$
vanishing. We apply formal GAGA over the completed local ring at
each exceptional image; no algebraization from completion along
an arbitrary divisor of the global surface is required. It algebraizes them
\citep[Tags 0885, 08MD, 08MF]{C_StacksBoundedReflexive}.

The same calculation with one bundle fixed controls the discrepancy
between $E$ and $\pi_*F$. Specifically, the extension identity off the
exceptional set gives
\[
 E=\pi_*F(BZ),\qquad
 E/\pi_*F\hookrightarrow\pi_*\bigl(F(BZ)/F\bigr).
\]
The dimensions of the cohomology of the fixed thickening on the right
are $O(r)$, by its fixed component filtration and the slope windows.
For $R^1\pi_*F$, formal functions and the second vanishing show that
the inverse system stabilizes at one fixed thickening $NZ$; its
$H^1$ is again $O(r)$. Also $R^2\pi_*F=0$, since the fibres have
dimension at most one. Thus
\[
 |\chi(Y,F)-\chi(S,E)|\le Cr.
\]
Frobenius scaling gives numerical $c_1(E)=0$ and
$\operatorname{ch}_2(E)=0$, independently of the period. Proper
Riemann--Roch with the rank-linear correction yields
$|\chi(S,E)|\le Cr$
\citep[Theorem 4.5(2)]{C_Langer2025}
\citep[Theorem 0.1]{C_Langer2026}.

Put $q=p^f$. The isomorphism on $Y\setminus Z$ between $F_Y^{q*}F$
and $F$ gives an actual line-bundle identity
\[
 (\det F)^{\otimes(q-1)}\simeq\mathcal O_Y(D_q)
\]
for an exceptional divisor $D_q$. Hence $c_1(F)$ belongs to the
exceptional numerical span. Its intersection with every $E_i$ is
bounded by a constant times $r$ from the slope windows. The negative
definite exceptional intersection matrix then bounds each coefficient
by $Cr$. This argument uses the Frobenius determinant identity; it
does not assume that numerical triviality downstairs on a nonprojective
surface already identifies the exceptional span upstairs.

The rank-independent one-step slope estimates for Frobenius on a fixed
smooth curve are those of
\citet[Theorem 2.2]{C_Sun2008} and
\citet[Proposition 1.3]{C_MehtaPauly2007}. Iteration bounds the slopes
of the restrictions of $F_Y^{q*}F$ to each exceptional component by
constants times $q$. The preceding pole-lowering argument therefore
extends the periodic isomorphism and its inverse with a pole bound
$B_qZ$, where $B_q\le Cq$. We obtain
\[
 (q-1)\mu_{\max,H}(F)\le B_q Z\cdot H
\]
and the same inequality for the dual. Since $q\ge2$, both slopes are
uniformly bounded. Choose a fixed smooth ample curve not contained in
the exceptional locus or the determinant-zero divisors. Restricting
the same maps gives the same estimates directly on this curve, without
a restriction theorem depending on the rank.

The numerical $c_1$ bounds and the Euler bound give finitely many
Hilbert polynomials for each fixed rank. Langer's boundedness theorem
applies on the fixed smooth projective $Y$
\citep[Theorem 4.2, printed p.~269]{C_LangerBounded2004}.
The full addendum explicitly retains this theorem
\citep{C_LangerAddendum2004}. We therefore choose one finite-type
parameter space with a bounded family of the relevant bundles. We
retain the open slope-window conditions; they can be expressed using
the proper images of finitely many Quot schemes of all coherent
quotients of fixed positive rank and degree below a threshold. Torsion
only raises the degree, and a uniform global-generation bound supplies
a lower bound on the degrees which need be tested. Taking only the
open torsion-free part of a Quot scheme would not justify closedness
of the bad locus.

Finally, bounded-pole Hom schemes on the proper $Y$ are of finite
presentation. Pairs of maps in the two directions, whose products equal
the fixed power of the section defining $Z$, encode ordinary
isomorphisms off $Z$. Their images in the pair of parameter spaces are
constructible. The same construction with $F_Y^*F$ gives the
one-step Frobenius relation: this source has fixed exceptional windows
by the one-step slope bound, so the required pole order is still fixed.
The relation is independent of the chosen minimal extensions because
both sides mean precisely the corresponding isomorphism of reflexive
sheaves downstairs. These assertions concern the represented
$k$-points; no universal algebraized family over a complete local base
is being asserted. This proves the final statement.
\end{proof}

\begin{theorem}[Proved: Fixed-rank growth of defining-characteristic representations]
\label{D:surface-fixed-rank-representation-growth}
Let $S$ be a proper scheme of dimension at most two over an algebraically
closed field of characteristic $p>0$, and let $T$ be finite of codimension
at least two in every component it meets. For a connected component
$U$ of $S\setminus T$ and each fixed $r$, there is $c_{U,r}$ such that
\[
 \#\bigl(\operatorname{Hom}_{\rm cont}
 (\pi_1(U),\operatorname{GL}_r(\mathbf F_{p^f}))/\mathrm{conj}\bigr)
 \le p^{c_{U,r}f}\qquad(f\ge1).
\]
\end{theorem}
\begin{proof}
First let $S$ be normal and work on its regular locus. Use the bounded
atlas in Proposition~\ref{D:global-periodic-reflexive-boundedness}.
Beke's effective quotient theorem for constructible equivalence
relations gives a constructible space $P_r$ whose $k$-points are the
represented reflexive isomorphism classes
\citep[Theorem 2]{Beke2011FixedRankB}. The Frobenius relation descends
to the graph of a partial constructible function $\Phi$ on $P_r$.
It may be undefined when a nonperiodic object leaves the bounded
family. Every periodic orbit is defined throughout and remains in
the family, by the same boundedness proposition.

A fixed periodic object admits at most one unit structure up to
isomorphism. To see this, set $A=H^0(S,\operatorname{End}E)$, a
finite-dimensional $k$-algebra. Its unit group is connected: its
semisimple quotient is a product of general linear groups, and the
kernel from the radical is unipotent. A chosen $q$-unit structure
defines a $q$-semilinear algebra endomorphism $\tau$ of $A$.
The differential of $\tau$ is zero. Every twisted-conjugation orbit
in $A^\times$ is therefore open, since the orbit map has invertible
differential. Connectedness makes it a single orbit. This is the
required Lang surjectivity and uniqueness.

Unit structures correspond to finite etale
$\operatorname{GL}_r(\mathbf F_q)$-frame torsors. The proved smallness
of punctured proper normal surfaces implies that there are finitely
many such classes for each $q$. The exact source is the combination
of ILO Expose XXI, Theorem 1.3, with its proper-composition
Lemma 3.2.1(ii), already used in Proposition~\ref{D:codim-two-small}.
Thus the $f$-cycles of $\Phi$ form a finite set of $k$-points.

Decompose the graph of $\Phi$ into finitely many affine locally closed
pieces, each defined by bounded-degree equations and inequations.
A length-$f$ word of pieces gives a cycle locus in $O_r(f)$ variables
with $O_r(f)$ equations of fixed degree. Replace inequations by inverse
variables. Its set of geometric points is finite, so its reduced
support is zero-dimensional. Successive hypersurface intersection
bounds its number of points exponentially in $f$. There are only
exponentially many words. Hence the total is at most $p^{c_rf}$.
This argument applies to a partial function; alternatively one can add
an absorbing point for the undefined domain.

Purity permits removal of finitely many regular points. Restriction
from a normal scheme to a dense open preserves connectedness of every
connected finite etale cover, so the regular-locus group surjects
onto $\pi_1(U)$. Finally finite normalization handles a nonnormal
surface. For each fixed collection of upstairs representation classes,
descent isomorphisms on the finitely many connected components of the
double fibre product have at most
$|\operatorname{GL}_r(\mathbf F_q)|^c\le q^{cr^2}$ choices; impose the
cocycle equation and apply effective finite-surjective descent
\citep[Expose IX, Theorem 4.7]{C_SGA1}. Proper lower-dimensional
components have finitely generated fundamental groups. The finite
product of these estimates proves the stated scope.
\end{proof}
